% !TeX program = pdflatex
\documentclass[aspectratio=169,xcolor={dvipsnames,table}]{beamer}

\input{../../_en_preambulo_beamer}
% Comandos y macros estadísticas compartidas para las presentaciones Beamer en inglés (Metropolis)

% Conjuntos numéricos básicos
\newcommand{\R}{\mathbb{R}}
\newcommand{\N}{\mathbb{N}}
\newcommand{\Q}{\mathbb{Q}}
\newcommand{\Z}{\mathbb{Z}}
\newcommand{\set}[1]{\left\{ #1 \right\}}
\newcommand{\abs}[1]{\left|#1\right|}
\newcommand{\norm}[1]{\left\| #1 \right\|}

% Comandos estadísticos y probabilísticos del curso
\newcommand{\Var}{\operatorname{Var}}
\newcommand{\cov}{\operatorname{Cov}}
\newcommand{\corr}{\rho}
\newcommand{\comb}[2]{\begin{pmatrix} #1 \\ #2 \end{pmatrix}}
\newcommand{\s}{\sigma}
\newcommand{\particion}{\mathfrak{P}}
\newcommand{\card}[1]{\#\left( #1 \right)}
\newcommand{\seq}[3]{\set{#1}_{#2}^{#3}}
\newcommand{\E}{\mathbb{E}}
\newcommand{\Prob}{\mathbb{P}}

% Entornos pedagógicos y cajas en inglés académico natural (soportando tanto nombres en inglés como en español)
\newenvironment{discussion}{%
  \begin{alertblock}{Discussion Question}%
}{%
  \end{alertblock}%
}
\newenvironment{discusion}{%
  \begin{alertblock}{Discussion Question}%
}{%
  \end{alertblock}%
}

\newenvironment{reflection}{%
  \begin{alertblock}{Classroom Reflection}%
}{%
  \end{alertblock}%
}
\newenvironment{reflexion}{%
  \begin{alertblock}{Classroom Reflection}%
}{%
  \end{alertblock}%
}

\newenvironment{paradox}{%
  \begin{alertblock}{Exploratory Paradox}%
}{%
  \end{alertblock}%
}
\newenvironment{paradoja}{%
  \begin{alertblock}{Exploratory Paradox}%
}{%
  \end{alertblock}%
}

\newenvironment{motivation}{%
  \begin{exampleblock}{Motivation: Data Science in Practice}%
}{%
  \end{exampleblock}%
}
\newenvironment{motivacion}{%
  \begin{exampleblock}{Motivation: Data Science in Practice}%
}{%
  \end{exampleblock}%
}

\newenvironment{quickchallenge}{%
  \begin{block}{Quick Team Challenge}%
}{%
  \end{block}%
}
\newenvironment{retoclase}{%
  \begin{block}{Quick Team Challenge}%
}{%
  \end{block}%
}


\title{Hypergeometric Distribution}
\subtitle{Section 03.06 --- Sampling Without Replacement and FPCF}
\author[J. Castillo Colmenares]{Juliho Castillo Colmenares}
\institute[Tec de Monterrey]{Tecnológico de Monterrey}
\date{\vspace{-1.5cm}}

\begin{document}

% SLIDE 01: Title Page
\begin{frame}[plain]
	\titlepage
\end{frame}

% SLIDE 02: Roadmap
\begin{frame}{Roadmap --- Chapter 03: Discrete Random Variables}
	\vspace{-0.15cm}
	\begin{block}{Curricular Structure of Chapter 03}
		\footnotesize
		\begin{enumerate}
			\itemsep=0.03cm
			\item \textcolor{gray}{Section 03.01: Discrete Probability Mass Functions (PMF and Support)}
			\item \textcolor{gray}{Section 03.02: Cumulative Distribution Functions (Discrete CDF)}
			\item \textcolor{gray}{Section 03.03: Mathematical Expectation, Variance, and Moments}
			\item \textcolor{gray}{Section 03.04: Bernoulli and Binomial Distributions}
			\item \textcolor{gray}{Section 03.05: Geometric and Negative Binomial Distributions}
			\item \textbf{\textcolor{TecRojo}{Section 03.06: Hypergeometric Distribution}} $\leftarrow$ \emph{Current focus}
		\end{enumerate}
	\end{block}
	\vspace{-0.2cm}
	\begin{alertblock}{Learning Objective}
		\scriptsize
		Master the modeling of random draws \emph{without replacement} from heterogeneous finite populations, quantifying variance reduction through the Finite Population Correction Factor (FPCF) and negative covariance.
	\end{alertblock}
\end{frame}

% SLIDE 03: Motivation and Intuition
\begin{frame}{Motivation --- From Independence to Sample Dependence}
	\vspace{-0.15cm}
	\begin{columns}[T]
		\begin{column}{0.48\textwidth}
			\begin{block}{Sampling With Replacement (Binomial)}
				\footnotesize
				\begin{itemize}
					\itemsep=0.05cm
					\item Each drawn element returns to the urn before the next draw.
					\item Success probability ($p$) remains **strictly constant** across draws.
					\item Trials are mutually **independent** ($P(A \cap B) = P(A)P(B)$).
					\item Theoretical variance maximum: $\sigma^2 = np(1-p)$.
				\end{itemize}
			\end{block}
		\end{column}
		\begin{column}{0.48\textwidth}
			\begin{alertblock}{Sampling Without Replacement (Hypergeometric)}
				\footnotesize
				\begin{itemize}
					\itemsep=0.05cm
					\item Drawn elements **do not return**; population shrinks ($N \to N-1$).
					\item Probability of the next attempt depends heavily on prior history.
					\item Trials exhibit **negative correlated dependence**.
					\item Variance reduced by the population structure.
				\end{itemize}
			\end{alertblock}
		\end{column}
	\end{columns}
\end{frame}

% SLIDE 04: Formal Definition
\begin{frame}{Formal Definition --- The Hypergeometric Probability Mass Function}
	\vspace{-0.2cm}
	\begin{block}{Population Structure ($N, K, n$)}
		\scriptsize
		In a finite population of size $N$, exactly $K$ elements are \emph{successes} and $N-K$ are \emph{failures}. Drawing a sample of size $n$ without replacement, the success count $X$ follows a **Hypergeometric Distribution**: $X \sim \text{Hyper}(N, K, n)$.
	\end{block}
	\vspace{-0.25cm}
	\begin{columns}[T]
		\begin{column}{0.48\textwidth}
			\begin{exampleblock}{Exact PMF Formula}
				\scriptsize
				Probability of observing $k$ successes:
				$P(X = k) = \frac{\comb{K}{k}\comb{N-K}{n-k}}{\comb{N}{n}}$.
				
				Where $\comb{K}{k}$ chooses successes, $\comb{N-K}{n-k}$ chooses failures, and $\comb{N}{n}$ is total sample space.
			\end{exampleblock}
		\end{column}
		\begin{column}{0.48\textwidth}
			\begin{alertblock}{SciPy Parameterization}
				\scriptsize
				In Python (\texttt{scipy.stats.hypergeom}):
				$\texttt{hypergeom(M, n, N)}$.
				
				Variable mapping:
				\texttt{M} $= N$ (Total population),
				\texttt{n} $= K$ (Total successes),
				\texttt{N} $= n$ (Sample size).
			\end{alertblock}
		\end{column}
	\end{columns}
\end{frame}

% SLIDE 05: Exact Support and Vandermonde
\begin{frame}{Exact Support and Vandermonde Normalization}
	\vspace{-0.15cm}
	\begin{columns}[T]
		\begin{column}{0.48\textwidth}
			\begin{block}{Support Interval $\mathcal{S}_X$}
				\scriptsize
				Combinatorial coefficients impose a dual physical constraint on $k$:
				\begin{itemize}
					\itemsep=0.04cm
					\item Available successes: $k \le \min(n, K)$.
					\item Available failures: $n-k \le N-K \implies k \ge n-(N-K)$.
				\end{itemize}
				The exact discrete support is:
				\begin{align*}
					\mathcal{S}_X = \set{\max(0, n-N+K), \dots, \min(n, K)}
				\end{align*}
			\end{block}
		\end{column}
		\begin{column}{0.48\textwidth}
			\begin{exampleblock}{Vandermonde Identity}
				\scriptsize
				To guarantee total probability sums to 1, we apply the combinatorial convolution identity with $a=K$, $b=N-K$, $c=n$:
				\begin{align*}
					\sum_{k} \comb{a}{k} \comb{b}{c-k} = \comb{a+b}{c}
				\end{align*}
				\vspace{-0.25cm}
				Dividing by $\comb{N}{n}$, we obtain:
				\begin{align*}
					\sum_{k \in \mathcal{S}_X} P(X=k) = 1
				\end{align*}
			\end{exampleblock}
		\end{column}
	\end{columns}
\end{frame}

% SLIDE 06: Expectation and Linearity
\begin{frame}{Mathematical Expectation --- Symmetry and Unconditional Linearity}
	\vspace{-0.15cm}
	\begin{block}{Decomposition via Indicator Variables}
		\footnotesize
		Let $I_i = 1$ if the $i$-th draw is a success ($0$ otherwise), so $X = \sum_{i=1}^n I_i$. By unconditional marginal symmetry, every element has an identical initial probability of being chosen at draw $i$: $P(I_i = 1) = \frac{K}{N}$.
	\end{block}
	\vspace{-0.2cm}
	\begin{columns}[T]
		\begin{column}{0.48\textwidth}
			\begin{exampleblock}{First Moment Derivation}
				\scriptsize
				By universal linearity of expectation:
				\begin{align*}
					\mathbb{E}[X] = \sum_{i=1}^n \mathbb{E}[I_i] = \sum_{i=1}^n \dfrac{K}{N} = n \dfrac{K}{N}
				\end{align*}
				Matches binomial expectation ($np$) exactly when setting $p = K/N$!
			\end{exampleblock}
		\end{column}
		\begin{column}{0.48\textwidth}
			\begin{alertblock}{Symmetry and Conditioning}
				\scriptsize
				Although draw 1 alters draw 2 ($P(I_2=1 \mid I_1=1) = \frac{K-1}{N-1}$), total probability preserves the marginal:
				\begin{align*}
					P(I_2=1) = \dfrac{K}{N}\dfrac{K-1}{N-1} + \dfrac{N-K}{N}\dfrac{K}{N-1} = \dfrac{K}{N}
				\end{align*}
			\end{alertblock}
		\end{column}
	\end{columns}
\end{frame}

% SLIDE 07: Variance and Negative Covariance
\begin{frame}{Negative Covariance and the FPCF Factor}
	\vspace{-0.15cm}
	\begin{columns}[T]
		\begin{column}{0.48\textwidth}
			\begin{block}{Sample Covariance ($i \neq j$)}
				\scriptsize
				Joint probability is $P(I_i=1, I_j=1) = \frac{K(K-1)}{N(N-1)}$. Its covariance yields:
				\begin{align*}
					\text{Cov}(I_i, I_j) &= \dfrac{K(K-1)}{N(N-1)} - \left(\dfrac{K}{N}\right)^2 \\
					&= -\dfrac{K(N-K)}{N^2(N-1)} < 0
				\end{align*}
				Extracting one success reduces the likelihood of extracting another later.
			\end{block}
		\end{column}
		\begin{column}{0.48\textwidth}
			\begin{exampleblock}{Variance and Population Correction}
				\scriptsize
				Summing individual variances plus $n(n-1)$ negative covariances:
				\begin{align*}
					\text{Var}(X) &= n\dfrac{K}{N}\left(1-\dfrac{K}{N}\right) \left[ \dfrac{N-n}{N-1} \right] \\
					&= np(1-p) \times \text{FPCF}
				\end{align*}
				\vspace{-0.25cm}
				The term $\text{FPCF} = \frac{N-n}{N-1} \le 1$ quantifies theoretical uncertainty reduction.
			\end{exampleblock}
		\end{column}
	\end{columns}
\end{frame}

% SLIDE 08: Theoretical Comparison Binomial vs Hypergeometric
\begin{frame}{Uncertainty Analysis --- Census vs. Sampling}
	\vspace{-0.15cm}
	\begin{block}{Limit Behavior of the FPCF ($\frac{N-n}{N-1}$)}
		\footnotesize
		The Finite Population Correction Factor acts as a geometric variance modulator driven by sample size $n$:
	\end{block}
	\vspace{-0.2cm}
	\begin{columns}[T]
		\begin{column}{0.48\textwidth}
			\begin{exampleblock}{Extreme Cases ($n=1$ and $n=N$)}
				\scriptsize
				\begin{itemize}
					\itemsep=0.04cm
					\item **Single Draw ($n=1$):** FPCF equals $\frac{N-1}{N-1} = 1$. Without and with replacement are identical for one attempt.
					\item **Full Census ($n=N$):** FPCF equals $0 \implies \text{Var}(X) = 0$. Measuring the whole population leaves zero random error ($X=K$ certainly).
				\end{itemize}
			\end{exampleblock}
		\end{column}
		\begin{column}{0.48\textwidth}
			\begin{alertblock}{Underdispersion vs Binomial}
				\scriptsize
				For a $40\%$ sample size ($n = 0.40 N$ with $N=100$):
				\begin{align*}
					\text{FPCF} = \dfrac{100-40}{100-1} = \dfrac{60}{99} \approx 0.606
				\end{align*}
				Sample variance is $39.4\%$ smaller than if we assumed sampling with replacement. The model is **underdispersed** ($\sigma^2 < \mu$)!
			\end{alertblock}
		\end{column}
	\end{columns}
\end{frame}

% SLIDE 09: Asymptotic Convergence
\begin{frame}{Asymptotic Convergence --- The $5\%$ Rule of Thumb}
	\vspace{-0.15cm}
	\begin{columns}[T]
		\begin{column}{0.48\textwidth}
			\begin{block}{Population Asymptotic Limit}
				\footnotesize
				When $N \to \infty$ and $K \to \infty$ holding proportion $p = \frac{K}{N}$ fixed, the FPCF converges to 1:
				\begin{align*}
					\lim_{N \to \infty} \dfrac{N-n}{N-1} = 1
				\end{align*}
				\vspace{-0.25cm}
				Drawing without replacement from an infinite population barely alters proportions, and $\text{Hyper}(N, K, n) \to \text{Bin}(n, p)$.
			\end{block}
		\end{column}
		\begin{column}{0.48\textwidth}
			\begin{alertblock}{Auditing Golden Rule ($\frac{n}{N} < 0.05$)}
				\footnotesize
				In quality engineering, computing massive factorials inside $\binom{N}{n}$ triggers numerical overflows.
				
				**Practical Rule:** If sample size represents less than **5\%** of the total population ($n/N < 0.05$), the hypergeometric can be approximated by the binomial with absolute error $< 10^{-3}$.
			\end{alertblock}
		\end{column}
	\end{columns}
\end{frame}

% SLIDE 10: Key Applications
\begin{frame}{Applications --- Quality Audits and Fisher's Exact Test}
	\vspace{-0.15cm}
	\begin{block}{Three Critical Domains of the Hypergeometric Distribution}
		\footnotesize
		\begin{enumerate}
			\itemsep=0.08cm
			\item **Quality Control (MIL-STD-105):** Attribute inspection in closed finite batches. From a batch of $N$ items containing $K$ defectives, a sample $n$ is audited and rejected if defects exceed limit $c$.
			\item **Forensic and Financial Auditing:** Evaluating finite transaction ledgers or invoice databases to quantify fiscal irregularities with exact mathematical rigor without overestimating variance.
			\item **Fisher's Exact Test in Biostatistics:** In $2 \times 2$ clinical trials with small sample sizes, conditioned on fixed marginal totals, the null distribution of event counts is **exactly hypergeometric**.
		\end{enumerate}
	\end{block}
\end{frame}

% SLIDE 11A: Classroom Exercise --- Tier 1: Fundamental (Statement)
\begin{frame}{Classroom Exercise --- Tier 1: Fundamental (1/2)}
	\footnotesize
	\begin{block}{Problem 3.6.2 --- Industrial Quality Batch Acceptance (Statement)}
		A receiving batch contains $N = 30$ electronic sensors, of which exactly $K = 6$ are defective ($20\%$). A quality engineer selects a random sample of $n = 5$ sensors \emph{without replacement}. The batch is accepted if at most 1 defective sensor is found ($X \le 1$). Compute the exact probability of accepting the batch.
	\end{block}

	\vspace{0.2cm}
	\pause
	\begin{alertblock}{Setting Up the Exact Model}
		The random variable $X$ represents the number of defective sensors in the sample. Since elements are drawn without replacement from a finite population, it follows a hypergeometric law:
		$$X \sim \text{Hyper}(N=30, \; K=6, \; n=5)$$ \pause
		The total acceptance probability corresponds to the sum for $X \le 1$:
		$$P(\text{Accept}) = P(X = 0) + P(X = 1)$$
	\end{alertblock}
\end{frame}

% SLIDE 11B: Classroom Exercise --- Tier 1: Fundamental (Solution)
\begin{frame}{Classroom Exercise --- Tier 1: Fundamental (2/2)}
	\scriptsize
	We compute the point probabilities using the combinatorial PMF formula: \pause

	\vspace{0.04cm}
	\begin{block}{Evaluation of Combinatorial Probabilities}
		\begin{itemize}\itemsep=0.04cm
			\item \textbf{Zero defectives ($X=0$):}
			$P(X=0) = \dfrac{\comb{6}{0}\comb{24}{5}}{\comb{30}{5}} = \dfrac{1 \times 42,504}{142,506} \approx 0.29826$ \pause
			\item \textbf{One defective ($X=1$):}
			$P(X=1) = \dfrac{\comb{6}{1}\comb{24}{4}}{\comb{30}{5}} = \dfrac{6 \times 10,626}{142,506} \approx 0.44739$
		\end{itemize}
	\end{block}

	\vspace{0.04cm}
	\pause
	\begin{alertblock}{Engineering Conclusion}
		Summing both success probabilities of the quality control:
		$$P(\text{Accept}) = 0.29826 + 0.44739 = \mathbf{0.74565} \; (74.57\%)$$
	\end{alertblock}
\end{frame}

% SLIDE 12A: Classroom Exercise --- Tier 2: Operational (Statement)
\begin{frame}{Classroom Exercise --- Tier 2: Operational (1/2)}
	\footnotesize
	\begin{block}{Problem 3.6.4 --- Variance Reduction via FPCF (Statement)}
		A shipping company has a fleet of $N = 50$ delivery trucks, of which $K = 10$ have brake system issues ($p = 0.20$). The inspector randomly selects $n = 20$ trucks for detailed inspection without replacement. Quantitatively compare the variance of trucks with issues under the hypergeometric model versus the binomial model.
	\end{block}

	\vspace{0.2cm}
	\pause
	\begin{alertblock}{Variance Under the Binomial Assumption (With Replacement)}
		If sampling were done with replacement (or if the population were infinite), the variance would follow independent Bernoulli trials:
		$$\text{Var}_{\text{Bin}}(X) = n p (1-p) = 20(0.20)(0.80) = \mathbf{3.20}$$
	\end{alertblock}
\end{frame}

% SLIDE 12B: Classroom Exercise --- Tier 2: Operational (Solution)
\begin{frame}{Classroom Exercise --- Tier 2: Operational (2/2)}
	\footnotesize
	We now incorporate the exact nature of sampling from a finite population without replacement: \pause

	\vspace{0.15cm}
	\begin{block}{Correction Factor and Exact Variance}
		\begin{itemize}\itemsep=0.15cm
			\item \textbf{Finite Population Correction Factor (FPCF):}
			$$\text{FPCF} = \dfrac{N-n}{N-1} = \dfrac{50-20}{50-1} = \dfrac{30}{49} \approx \mathbf{0.61224}$$ \pause
			\item \textbf{Exact Hypergeometric Variance:}
			$$\text{Var}_{\text{Hyp}}(X) = \text{Var}_{\text{Bin}}(X) \times \text{FPCF} = 3.20 \times 0.61224 = \mathbf{1.95918}$$
		\end{itemize}
	\end{block}

	\vspace{0.15cm}
	\pause
	\begin{alertblock}{Uncertainty Reduction Analysis}
		Sampling without replacement reduces estimation uncertainty by exactly $\mathbf{38.78\%}$ compared to the conventional binomial model, by eliminating double-counting of trucks.
	\end{alertblock}
\end{frame}

% SLIDE 13A: Classroom Exercise --- Tier 3: Analytical (Statement)
\begin{frame}{Classroom Exercise --- Tier 3: Analytical (1/2)}
	\footnotesize
	\begin{block}{Problem 3.6.7 --- Proof of the Factorial Moment (Statement)}
		Starting from the fundamental combinatorial identity $k \comb{K}{k} = K \comb{K-1}{k-1}$, algebraically prove that the first factorial moment of a hypergeometric variable is exactly $\mathbb{E}[X] = n \dfrac{K}{N}$.
	\end{block}

	\vspace{0.1cm}
	\pause
	\begin{alertblock}{Step 1: Expanding the Expectation Definition}
		\scriptsize
		We start from the weighted sum of probabilities over the support $k > 0$:
		$$\mathbb{E}[X] = \sum_{k=1}^{\min(n,K)} k \dfrac{\comb{K}{k}\comb{N-K}{n-k}}{\comb{N}{n}}$$ \pause
		We apply absorption in the numerator $k\comb{K}{k} = K\comb{K-1}{k-1}$ and in the denominator $n\comb{N}{n} = N\comb{N-1}{n-1} \implies \comb{N}{n} = \frac{N}{n}\comb{N-1}{n-1}$.
	\end{alertblock}
\end{frame}

% SLIDE 13B: Classroom Exercise --- Tier 3: Analytical (Solution)
\begin{frame}{Classroom Exercise --- Tier 3: Analytical (2/2)}
	\scriptsize
	Substituting both combinatorial identities into the expectation expression, we obtain: \pause

	\vspace{-0.15cm}
	\begin{block}{Step 2: Extraction and Vandermonde Sum}
		\scriptsize
		$\mathbb{E}[X] = \dfrac{n K}{N} \sum_{k=1}^{\min(n,K)} \dfrac{\comb{K-1}{k-1}\comb{N-K}{n-k}}{\comb{N-1}{n-1}}$ \pause \\
		\vspace{0.02cm}
		With index $j = k-1$, the sum spans all samples of size $n-1$: \\
		\vspace{0.02cm}
		$\mathbb{E}[X] = \dfrac{n K}{N} \sum_{j=0}^{\min(n-1,K-1)} \dfrac{\comb{K-1}{j}\comb{(N-1)-(K-1)}{(n-1)-j}}{\comb{N-1}{n-1}}$
	\end{block}

	\vspace{-0.15cm}
	\pause
	\begin{alertblock}{Conclusion via Unit Normalization}
		\scriptsize
		The sum is the total probability of $\text{Hyper}(N-1, K-1, n-1)$ and equals $1 \implies \mathbb{E}[X] = n \dfrac{K}{N} \cdot 1 = np \quad \blacksquare$
	\end{alertblock}
\end{frame}

% SLIDE 14A: Classroom Exercise --- Tier 4: Challenging (Statement)
\begin{frame}{Classroom Exercise --- Tier 4: Challenging (1/2)}
	\footnotesize
	\begin{block}{Problem 3.6.10 --- Fisher's Exact Test in Biomedical Trials}
		A cardiology clinical trial evaluates a new antithrombotic therapy with $N = 25$ severe patients. $K = 12$ patients are randomly assigned to the active drug and the remaining $13$ to placebo. At the end of the trial, exactly $n = 8$ patients developed deep thrombosis. If the drug were ineffective ($H_0$), the 8 cases of thrombosis would be distributed at random.
	\end{block}

	\vspace{0.2cm}
	\pause
	\begin{alertblock}{Setting Up the Hypothesis and Exact Inference}
		Let $X$ be the number of patients in the active group who developed thrombosis. Under the null hypothesis $H_0$, the variable is strictly hypergeometric:
		$$X \sim \text{Hyper}(N=25, \; K=12, \; n=8)$$ \pause
		In our clinical trial, we observe only $X = 1$ thrombosis case in the active group (and the other 7 in the placebo group). Let us evaluate the exact statistical significance.
	\end{alertblock}
\end{frame}

% SLIDE 14B: Classroom Exercise --- Tier 4: Challenging (Solution)
\begin{frame}{Classroom Exercise --- Tier 4: Challenging (2/2)}
	\scriptsize
	To determine significance at level $\alpha = 0.05$, we compute the exact one-tailed $p$-value ($P(X \le 1)$ under $H_0$): \pause

	\vspace{0.04cm}
	\begin{block}{Computing the $p$-value via the Hypergeometric Law}
		\scriptsize
		$p\text{-value} = P(X = 0) + P(X = 1) = \dfrac{\comb{12}{0}\comb{13}{8}}{\comb{25}{8}} + \dfrac{\comb{12}{1}\comb{13}{7}}{\comb{25}{8}}$ \pause
		\begin{itemize}\itemsep=0.01cm
			\item $P(X=0) = \frac{1 \times 1,287}{1,081,575} \approx 0.00119$ \quad | \quad $P(X=1) = \frac{12 \times 1,716}{1,081,575} \approx 0.01904$
		\end{itemize}
	\end{block}

	\vspace{0.04cm}
	\pause
	\begin{alertblock}{Clinical Decision ($p\text{-value} < \alpha$)}
		\scriptsize
		Since $p\text{-value} = 0.00119 + 0.01904 = \mathbf{0.02023}$ ($2.02\%$) is less than $\alpha = 0.05$, we reject $H_0$. We conclude that the reduction in thrombosis under the active therapy is clinically significant.
	\end{alertblock}
\end{frame}

% SLIDE 17: Python Lab I --- PMF and Support
\begin{frame}[fragile]{Python Lab --- Vectorized PMF/Support Validation (`scipy.stats`)}
	\vspace{-0.2cm}
	\begin{block}{Implementation in \texttt{03.06\_hypergeometric.py} (Block 1)}
		\lstinputlisting[language=Python, firstline=16, lastline=33, basicstyle=\fontsize{6.2pt}{7.2pt}\ttfamily]{../../code/03_variables_aleatorias_discretas/03.06_hypergeometric.py}
	\end{block}
\end{frame}

% SLIDE 18: Python Lab II --- FPCF and Monte Carlo
\begin{frame}[fragile]{Python Lab --- Negative Covariance and FPCF via Monte Carlo ($N=250,000$)}
	\vspace{-0.25cm}
	\begin{block}{Implementation in \texttt{03.06\_hypergeometric.py} (Block 2)}
		\lstinputlisting[language=Python, firstline=37, lastline=58, basicstyle=\fontsize{5.5pt}{6.5pt}\ttfamily]{../../code/03_variables_aleatorias_discretas/03.06_hypergeometric.py}
	\end{block}
\end{frame}

% SLIDE 19: Python Lab III --- Fisher's Test and Asymptotics
\begin{frame}[fragile]{Python Lab --- Fisher's Exact Test and Asymptotic Convergence}
	\vspace{-0.2cm}
	\begin{block}{Implementation in \texttt{03.06\_hypergeometric.py} (Block 3)}
		\lstinputlisting[language=Python, firstline=62, lastline=78, basicstyle=\fontsize{6.2pt}{7.2pt}\ttfamily]{../../code/03_variables_aleatorias_discretas/03.06_hypergeometric.py}
	\end{block}
\end{frame}

% SLIDE 20: Closing and Didactic Bridge
\begin{frame}{Closing --- Stochastic Synthesis and Didactic Bridge}
	\vspace{-0.15cm}
	\begin{columns}[T]
		\begin{column}{0.48\textwidth}
			\begin{block}{Stochastic Synthesis of Hypergeometric}
				\footnotesize
				\begin{itemize}
					\itemsep=0.08cm
					\item **Sampling Without Replacement ($N, K, n$):** Exact combinatorial modeling in finite populations.
					\item **Underdispersion and FPCF:** $\text{Var}(X) = np(1-p)\left(\frac{N-n}{N-1}\right) < np(1-p)$ due to negative covariance across draws.
					\item **Binomial Convergence:** When $\frac{n}{N} < 0.05$, sampling without replacement approximates replacement.
				\end{itemize}
			\end{block}
		\end{column}
		\begin{column}{0.48\textwidth}
			\begin{alertblock}{Bridge to Section 03.07}
				\footnotesize
				We have finalized discrete counting models under fixed success probability $p$ or finite populations with known $N, K$.
				
				What happens in continuous phenomena where discrete events occur randomly across **time or physical space**, governed by a constant average rate $\lambda$?
				
				$\to$ This leads naturally to the **Poisson Distribution**.
			\end{alertblock}
		\end{column}
	\end{columns}
\end{frame}

\end{document}
